% =========================================================
%  ppt/12-aset-animasi-ui.tex
%  Materi 12 — Impor Aset Visual, Animasi, dan UI Dasar
%  Bab XIV buku ajar
% =========================================================
\documentclass[aspectratio=169,10pt]{beamer}
% =========================================================
%  ppt/config/packages.tex
%  Paket bersama untuk seluruh presentasi kuliah
%  Pengantar Pemrograman Game
% =========================================================

\usepackage[T1]{fontenc}
\usepackage[utf8]{inputenc}
\usepackage[indonesian]{babel}
\usepackage{lmodern}
\usepackage{textcomp}
\usepackage{amsmath}
\usepackage{amssymb}
\usepackage{graphicx}
\usepackage{booktabs}
\usepackage{tabularx}
\usepackage{array}
\usepackage{multicol}
\usepackage{listings}
\usepackage{upquote}
\usepackage{adjustbox}
\usepackage{tcolorbox}
\usepackage{tikz}
\usetikzlibrary{
  arrows.meta,
  positioning,
  shapes.geometric,
  shapes.misc,
  calc,
  fit,
  backgrounds,
  matrix,
  chains,
  decorations.pathreplacing,
  patterns
}

\tcbuselibrary{skins}

% Jangan memuat geometry: ukuran halaman diatur oleh beamer.

% =========================================================
%  ppt/config/theme.tex
%  Satu tema Beamer untuk keempat belas presentasi
%  Pengantar Pemrograman Game
% =========================================================

% --------------------------------------------------------- warna
\definecolor{ubNavy}{HTML}{1F3864}
\definecolor{ubAccent}{HTML}{2E75B6}
\definecolor{ubHead}{HTML}{D9E2F3}
\definecolor{ubSoft}{HTML}{F2F6FC}
\definecolor{ubLine}{HTML}{BFBFBF}
\definecolor{ubText}{HTML}{1F2937}
\definecolor{ubMuted}{HTML}{5A6678}
\definecolor{ubGreen}{HTML}{1E7B45}
\definecolor{ubGreenBg}{HTML}{F1F8F2}
\definecolor{ubRed}{HTML}{B3261E}
\definecolor{ubRedBg}{HTML}{FDF2F1}
\definecolor{ubWarn}{HTML}{A35A05}
\definecolor{ubWarnBg}{HTML}{FFF8EC}
\definecolor{ubGray}{HTML}{F6F6F6}
\definecolor{ubCodeBg}{HTML}{10233F}
\definecolor{ubCodeText}{HTML}{E8EEF7}
\definecolor{ubCodeKey}{HTML}{8FBCF0}
\definecolor{ubCodeCom}{HTML}{8FB37A}
\definecolor{ubCodeStr}{HTML}{D9A46B}
\definecolor{ubCodeNum}{HTML}{7E99BE}

% Nama warna yang dipakai gambar TikZ buku ajar.
\definecolor{ubPrimary}{HTML}{1F3864}
\definecolor{ubAlt}{HTML}{F2F6FC}

% --------------------------------------------------------- dasar
\usetheme{default}
\usefonttheme{professionalfonts}
\setbeamertemplate{navigation symbols}{}
\setbeamerfont{frametitle}{size=\Large,series=\bfseries}
\setbeamerfont{framesubtitle}{size=\small}
\setbeamerfont{footline}{size=\tiny}
\setbeamerfont{block title}{size=\small,series=\bfseries}
\setbeamerfont{block body}{size=\small}
\setbeamerfont{itemize/enumerate body}{size=\small}
\setbeamerfont{itemize/enumerate subbody}{size=\footnotesize}

\setbeamersize{text margin left=0.60cm,text margin right=0.60cm}
\uchyph=0

\setbeamercolor{normal text}{fg=ubText,bg=white}
\setbeamercolor{structure}{fg=ubAccent}
\setbeamercolor{frametitle}{fg=ubNavy,bg=white}
\setbeamercolor{title}{fg=ubNavy}
\setbeamercolor{footline}{fg=ubMuted,bg=white}
\setbeamercolor{itemize item}{fg=ubAccent}
\setbeamercolor{itemize subitem}{fg=ubAccent}
\setbeamercolor{block title}{fg=white,bg=ubNavy}
\setbeamercolor{block body}{fg=ubText,bg=ubSoft}
\setbeamercolor{block title alerted}{fg=white,bg=ubWarn}
\setbeamercolor{block body alerted}{fg=ubText,bg=ubWarnBg}
\setbeamercolor{block title example}{fg=white,bg=ubGreen}
\setbeamercolor{block body example}{fg=ubText,bg=ubGreenBg}
\setbeamercolor{caption name}{fg=ubMuted}

\setbeamertemplate{blocks}[rounded][shadow=false]
\setbeamertemplate{itemize item}{\raisebox{0.10em}{\color{ubAccent}\scriptsize$\blacksquare$}}
\setbeamertemplate{itemize subitem}{\raisebox{0.06em}{\color{ubAccent}\tiny$\blacktriangleright$}}
\setbeamertemplate{enumerate item}{%
  \makebox[1.3em][l]{\color{ubAccent}\insertenumlabel.}}

\addtobeamertemplate{itemize/enumerate body begin}{%
  \setlength{\itemsep}{0.42em}\setlength{\parskip}{0pt}}{}

% --------------------------------------------------------- judul frame
\setbeamertemplate{frametitle}[default][left]

% --------------------------------------------------------- kaki
\newcommand{\materino}{00}
\newcommand{\materifoot}{Pengantar Pemrograman Game \textperiodcentered\ Materi \materino}

\setbeamertemplate{footline}{%
  \begin{beamercolorbox}[wd=\paperwidth,ht=0.015cm,dp=0cm]{footline}%
    \color{ubHead}\rule{\paperwidth}{0.4pt}%
  \end{beamercolorbox}%
  \begin{beamercolorbox}[wd=\paperwidth,ht=0.40cm,dp=0.14cm,%
                         leftskip=0.60cm,rightskip=0.60cm]{footline}%
    \usebeamerfont{footline}\color{ubMuted}\materifoot\hfill\insertframenumber%
  \end{beamercolorbox}%
}

% --------------------------------------------------------- makro buku ajar
\providecommand{\code}[1]{\texttt{\detokenize{#1}}}
\providecommand{\istilah}[1]{\emph{#1}}
\providecommand{\CPMK}[1]{CPMK\,#1}
\providecommand{\CPL}[1]{CPL\,#1}
\providecommand{\SUBCPMK}[1]{Sub-CPMK\,#1}

% --------------------------------------------------------- gaya kode
\lstdefinelanguage{CSharp}{
  morekeywords={abstract,as,base,bool,break,byte,case,catch,char,class,const,
    continue,decimal,default,delegate,do,double,else,enum,event,explicit,
    extern,false,finally,fixed,float,for,foreach,goto,if,implicit,in,int,
    interface,internal,is,lock,long,namespace,new,null,object,operator,out,
    override,params,private,protected,public,readonly,ref,return,sbyte,
    sealed,short,sizeof,stackalloc,static,string,struct,switch,this,throw,
    true,try,typeof,uint,ulong,unchecked,unsafe,ushort,using,virtual,void,
    volatile,while,var,get,set},
  sensitive=true,
  morecomment=[l]{//},
  morecomment=[s]{/*}{*/},
  morestring=[b]",
}

\lstdefinestyle{kode}{
  basicstyle=\ttfamily\footnotesize\color{ubCodeText},
  keywordstyle=\color{ubCodeKey}\bfseries,
  commentstyle=\color{ubCodeCom}\itshape,
  stringstyle=\color{ubCodeStr},
  numberstyle=\tiny\color{ubCodeNum},
  numbers=left,
  numbersep=7pt,
  backgroundcolor=\color{ubCodeBg},
  frame=none,
  showstringspaces=false,
  breaklines=false,
  tabsize=4,
  columns=fullflexible,
  keepspaces=true,
  xleftmargin=12pt,
  xrightmargin=2pt,
  aboveskip=0pt,
  belowskip=0pt,
  upquote=true,
}
\lstset{style=kode}

\newcommand{\kodejelas}[1]{%
  \vspace{0.16cm}%
  {\footnotesize\color{ubMuted}#1\par}%
}
\newcommand{\sumber}[1]{%
  \vspace{0.08cm}%
  {\scriptsize\color{ubMuted}\itshape #1\par}%
}
\newcommand{\awalan}[1]{%
  {\small\color{ubMuted}#1\par}\vspace{0.18cm}%
}

% --------------------------------------------------------- kotak
\newtcolorbox{kotak}[1][]{enhanced, colback=ubSoft, colframe=ubAccent,
  boxrule=0pt, leftrule=2.2pt, arc=1pt, left=6pt, right=6pt, top=4pt,
  bottom=4pt, fontupper=\footnotesize, #1}
\newtcolorbox{peringatan}[1][]{enhanced, colback=ubWarnBg, colframe=ubWarn,
  boxrule=0pt, leftrule=2.2pt, arc=1pt, left=6pt, right=6pt, top=4pt,
  bottom=4pt, fontupper=\footnotesize, #1}
\newtcolorbox{kiat}[1][]{enhanced, colback=ubGreenBg, colframe=ubGreen,
  boxrule=0pt, leftrule=2.2pt, arc=1pt, left=6pt, right=6pt, top=4pt,
  bottom=4pt, fontupper=\footnotesize, #1}

% --------------------------------------------------------- sampul
% \deckcover{nomor}{minggu}{judul}{sub-cpmk}{teknologi}{sumber}{blok}
\newcommand{\deckcover}[7]{%
\begin{frame}[plain]
\begin{tikzpicture}[remember picture, overlay]
  \fill[ubNavy] (current page.south west) rectangle ([xshift=6.00cm]current page.north west);
  \fill[ubAccent] ([xshift=6.00cm]current page.south west) rectangle ([xshift=6.08cm]current page.north west);

  \node[anchor=north west, text width=4.72cm, align=left, inner sep=0pt]
    at ($(current page.north west)+(0.66cm,-0.60cm)$) {%
      {\color{white}\bfseries\fontsize{10.2}{12.2}\selectfont
       Universitas Bale Bandung}\\[4pt]
      {\color{ubHead}\fontsize{7.6}{9.4}\selectfont
       Fakultas Teknologi Informasi}\\[1pt]
      {\color{ubHead}\fontsize{7.6}{9.4}\selectfont
       Program Studi Teknik Informatika}\\[7pt]
      {\color{ubAccent}\rule{4.0cm}{0.5pt}}%
    };

  \node[anchor=north west, text width=4.72cm, align=left, inner sep=0pt]
    at ($(current page.north west)+(0.66cm,-2.72cm)$) {%
      {\color{ubAccent}\fontsize{6.4}{7.8}\selectfont MATA KULIAH}\\[3pt]
      {\color{white}\bfseries\fontsize{9.8}{12.0}\selectfont
       Pengantar Pemrograman Game}\\[3pt]
      {\color{ubHead}\fontsize{7.6}{9.4}\selectfont
       MWCA009 \textperiodcentered\ 3 SKS}\\[8pt]
      {\color{ubHead}\fontsize{7.6}{9.4}\selectfont
       Semester Ganjil 2026/2027}%
    };

  \node[anchor=south west, text width=4.72cm, align=left, inner sep=0pt]
    at ($(current page.south west)+(0.66cm,0.52cm)$) {%
      {\color{ubAccent}\fontsize{6.4}{7.8}\selectfont DOSEN PENGAMPU}\\[3pt]
      {\color{white}\fontsize{8.4}{10.4}\selectfont
       Cecep Suwanda, S.Si., M.Kom.}%
    };

  \node[anchor=north west, text width=8.95cm, align=left, inner sep=0pt]
    at ($(current page.north east)+(-9.40cm,-0.92cm)$) {%
      {\color{ubAccent}\bfseries\fontsize{8.2}{10.2}\selectfont
       MATERI #1\ \ \textperiodcentered\ \ MINGGU #2}\\[8pt]
      {\color{ubNavy}\bfseries\fontsize{19.4}{23.4}\selectfont
       \hyphenpenalty=10000 \exhyphenpenalty=10000 \raggedright #3}%
    };

  \node[anchor=north west, text width=8.95cm, align=left, inner sep=0pt]
    at ($(current page.north east)+(-9.40cm,-4.30cm)$) {%
      \renewcommand{\arraystretch}{1.30}%
      \begin{tabular}{@{}>{\color{ubMuted}\fontsize{7.3}{9.0}\selectfont}p{1.85cm}
                        >{\color{ubText}\fontsize{8.1}{10.0}\selectfont}p{6.7cm}@{}}
        Sub-CPMK      & #4 \\
        Teknologi     & #5 \\
        Sumber materi & #6 \\
        Blok          & #7 \\
      \end{tabular}%
    };
\end{tikzpicture}
\end{frame}%
}

% --------------------------------------------------------- pembatas bagian
\newcommand{\sectionframe}[3]{%
\begin{frame}[plain]
\begin{tikzpicture}[remember picture, overlay]
  \fill[ubNavy] (current page.south west) rectangle (current page.north east);
  \fill[ubAccent] (current page.north west) rectangle ([xshift=0.10cm]current page.south west);
  \node[anchor=west, inner sep=0pt]
    at ($(current page.west)+(1.00cm,1.45cm)$)
    {\color{ubNavy!52!white}\bfseries\fontsize{58}{58}\selectfont #1};
  \node[anchor=west, inner sep=0pt, text width=11.6cm, align=left]
    at ($(current page.west)+(1.05cm,-0.55cm)$) {%
      {\color{ubAccent}\bfseries\fontsize{8.0}{10.0}\selectfont BAGIAN}\\[3pt]
      {\color{white}\bfseries\fontsize{20.5}{24.5}\selectfont #2}\\[5pt]
      {\color{ubHead}\small #3}%
    };
  \node[anchor=south east, inner sep=0pt]
    at ($(current page.south east)+(-0.60cm,0.32cm)$)
    {\color{ubHead}\fontsize{7.4}{9}\selectfont #1 / 6
     \ \ \textperiodcentered\ \ \insertframenumber};
\end{tikzpicture}
\end{frame}%
}

% --------------------------------------------------------- peta pertemuan
\newcommand{\blokpeta}[6]{%
  \ifnum#2=#3
    \node[rounded corners=3pt, fill=ubNavy, draw=ubNavy, line width=0.5pt,
          text width=3.02cm, minimum height=1.62cm, align=left, inner sep=5pt]
      at (#1,0) {%
        {\color{ubAccent}\bfseries\fontsize{6.8}{8}\selectfont #4}\\[1pt]
        {\color{white}\bfseries\fontsize{7.6}{9}\selectfont #5}\\[2pt]
        {\color{ubHead}\fontsize{6.4}{7.6}\selectfont #6}};
  \else
    \node[rounded corners=3pt, fill=ubSoft, draw=ubHead, line width=0.5pt,
          text width=3.02cm, minimum height=1.62cm, align=left, inner sep=5pt]
      at (#1,0) {%
        {\color{ubAccent}\bfseries\fontsize{6.8}{8}\selectfont #4}\\[1pt]
        {\color{ubNavy}\bfseries\fontsize{7.6}{9}\selectfont #5}\\[2pt]
        {\color{ubMuted}\fontsize{6.4}{7.6}\selectfont #6}};
  \fi
}

% \petapertemuan{blok aktif}{sebelumnya}{sesudahnya}{keluaran}
% \petapertemuan{blok aktif}{sebelumnya}{pertemuan ini}{sesudahnya}{keluaran}
\newcommand{\petapertemuan}[5]{%
\begin{frame}{Peta Pertemuan}
{\scriptsize\color{ubMuted}Kedudukan materi dalam empat blok perkuliahan}\par
\vspace{0.10cm}
\centering
\begin{tikzpicture}
  \blokpeta{-5.58}{1}{#1}{BLOK I}{Materi 1--5}{Game loop, vektor,\\fisika, collision}
  \blokpeta{-1.86}{2}{#1}{BLOK II}{Materi 6--9}{AI NPC: chase, FSM,\\pathfinding A*}
  \blokpeta{1.86}{3}{#1}{BLOK III}{Materi 10--11}{Unity: scene,\\level, physics}
  \blokpeta{5.58}{4}{#1}{BLOK IV}{Materi 12--14}{Aset visual, audio,\\UI, playtest}
\end{tikzpicture}
\par\vspace{0.30cm}
\begin{columns}[T,onlytextwidth]
  \begin{column}{0.315\textwidth}
    \begin{beamercolorbox}[wd=\linewidth,sep=4pt,rounded=true]{block body}
      {\color{ubMuted}\fontsize{6.6}{8}\selectfont SEBELUMNYA}\\[2pt]
      {\footnotesize #2}
    \end{beamercolorbox}
  \end{column}
  \begin{column}{0.315\textwidth}
    \begin{beamercolorbox}[wd=\linewidth,sep=4pt,rounded=true]{block title}
      {\color{ubHead}\fontsize{6.6}{8}\selectfont PERTEMUAN INI}\\[2pt]
      {\footnotesize #3}
    \end{beamercolorbox}
  \end{column}
  \begin{column}{0.315\textwidth}
    \begin{beamercolorbox}[wd=\linewidth,sep=4pt,rounded=true]{block body}
      {\color{ubMuted}\fontsize{6.6}{8}\selectfont SESUDAHNYA}\\[2pt]
      {\footnotesize #4}
    \end{beamercolorbox}
  \end{column}
\end{columns}
\par\vspace{0.26cm}
{\footnotesize\color{ubText}\textbf{Keluaran pertemuan.} #5\par}
\end{frame}%
}

% --------------------------------------------------------- penutup
\newcommand{\penutup}[2]{%
\begin{frame}[plain]
\begin{tikzpicture}[remember picture, overlay]
  \fill[ubNavy] (current page.south west) rectangle (current page.north east);
  \fill[ubAccent] (current page.north west) rectangle ([xshift=0.10cm]current page.south west);

  \node[anchor=north west, text width=13.6cm, align=left, inner sep=0pt]
    at ($(current page.north west)+(1.00cm,-0.95cm)$) {%
      {\color{ubAccent}\bfseries\fontsize{8.2}{10}\selectfont ASESMEN MINGGU INI}\\[6pt]
      {\color{white}\bfseries\fontsize{24}{28}\selectfont #1}%
    };

  \node[anchor=north west, text width=13.6cm, align=left, inner sep=0pt]
    at ($(current page.north west)+(1.00cm,-4.55cm)$) {%
      {\color{ubNavy!0!ubHead}\rule{13.4cm}{0.4pt}}\\[7pt]
      {\color{ubAccent}\bfseries\fontsize{7.4}{9}\selectfont
       PENGAIT KE MATERI BERIKUTNYA}\\[4pt]
      {\color{ubHead}\fontsize{10.6}{13.2}\selectfont #2}%
    };

  \node[anchor=south east, inner sep=0pt]
    at ($(current page.south east)+(-0.60cm,0.32cm)$)
    {\color{ubHead}\fontsize{7.4}{9}\selectfont \insertframenumber};
\end{tikzpicture}
\end{frame}%
}



\renewcommand{\materino}{12}

\begin{document}

\deckcover{12}{13}
  {Impor Aset Visual, Animasi, dan UI Dasar}
  {4.1}
  {Unity}
  {Bab XIV buku ajar}
  {Blok IV \textperiodcentered\ Materi 12--14}

% --------------------------------------------------------- capaian
\begin{frame}{Capaian Pertemuan}
\setbeamerfont{itemize/enumerate body}{size=\scriptsize}
\awalan{Sub-CPMK dan indikator dari RPS, serta tujuan Bab XIV buku ajar.}
\begin{columns}[T,onlytextwidth]
  \begin{column}{0.52\textwidth}
    \begin{beamercolorbox}[wd=\linewidth,sep=5pt,rounded=true]{block title}
      {\fontsize{6.8}{8}\selectfont\color{ubHead}SUB-CPMK 4.1 \textperiodcentered\ RANAH C3}\\[2pt]
      {\footnotesize Mengimpor dan mengonfigurasi aset visual
       (\emph{sprite}/model, animasi, UI) ke dalam \emph{scene}.}
    \end{beamercolorbox}
    \vspace{0.14cm}

    {\fontsize{7}{8.5}\selectfont\bfseries\color{ubAccent}INDIKATOR}\par
    \vspace{0.06cm}
    \begin{itemize}
      \item Mengimpor aset dan mengatur \emph{importer setting} dengan tepat.
      \item Menyiapkan \emph{Animator} beserta transisi animasi, misalnya
            diam dan berlari.
      \item Menampilkan elemen UI (skor dan nyawa) yang berfungsi pada
            \emph{scene}.
    \end{itemize}
    \vspace{0.10cm}
    \begin{kotak}
      Bab ini membuka CPMK 4. Yang dibahas bukan lagi mekanisme permainan,
      melainkan penyajiannya kepada pemain.
    \end{kotak}
  \end{column}
  \begin{column}{0.45\textwidth}
    {\fontsize{7}{8.5}\selectfont\bfseries\color{ubAccent}TUJUAN BAB XIV}\par
    \vspace{0.08cm}
    \begin{itemize}
      \item Menjelaskan jalan pengolahan aset, dari berkas gambar menjadi
            benda pada \emph{scene}.
      \item Mengatur empat \emph{importer setting} yang paling berpengaruh.
      \item Menjelaskan mengapa berkas aset tidak boleh diubah di luar Unity.
      \item Menyusun \emph{Animator} beserta keadaan dan transisinya.
      \item Menghubungkan keadaan pemain dengan animasi melalui parameter.
      \item Menyusun \emph{Canvas} yang menampilkan skor dan nyawa.
    \end{itemize}
  \end{column}
\end{columns}
\end{frame}

% --------------------------------------------------------- peta
\petapertemuan{4}
  {Materi 11 \textperiodcentered\ Membangun Level, \emph{Physics}, dan Kontrol
   Pemain}
  {Materi 12 \textperiodcentered\ Impor Aset Visual, Animasi, dan UI Dasar}
  {Materi 13 \textperiodcentered\ Integrasi Audio dan \emph{Trigger} Interaksi}
  {\emph{Scene} dengan aset visual sungguhan: animasi diam dan berlari
   berganti menurut gerak pemain, \emph{Canvas} menampilkan skor dan nyawa,
   dan \emph{placeholder} \emph{greybox} telah diganti.}

% ========================================================= bagian 1
\section{Jalan Pengolahan Aset}
\sectionframe{01}{Jalan Pengolahan Aset}
  {Dari berkas gambar menjadi benda pada scene}

\begin{frame}{Tiga Langkah dari Berkas menjadi Benda}
\awalan{Perbedaan besar antara Unity dan Pygame terletak pada saat aset
dikerjakan.}
\begin{enumerate}
  \item \textbf{Berkas disalin} ke dalam folder \texttt{Assets}. Unity
        memantau folder itu, sehingga berkas yang baru muncul langsung diolah.
  \item \textbf{Berkas diolah} menurut \emph{importer setting} yang berlaku
        baginya, lalu hasilnya disimpan pada folder \texttt{Library}. Pengolahan
        hanya dilakukan sekali untuk setiap pengaturan, bukan setiap kali
        \emph{scene} dijalankan.
  \item \textbf{Berkas dipakai} dengan menyeretnya dari \emph{Project window}
        ke \emph{Scene} atau ke komponen \emph{Sprite Renderer}.
\end{enumerate}
\vspace{0.12cm}
\begin{kotak}
Pada Pygame gambar dimuat ketika program dijalankan, sehingga program harus
disunting untuk mengubahnya. Pada Unity aset diolah ketika diimpor, dan
hasilnya disimpan.
\end{kotak}
\end{frame}

\begin{frame}{Empat Pengaturan Impor yang Paling Berpengaruh}
{\scriptsize
\renewcommand{\arraystretch}{1.14}
\begin{tabularx}{\textwidth}{@{}>{\raggedright\arraybackslash}p{2.9cm}
    >{\raggedright\arraybackslash}p{3.4cm}
    >{\raggedright\arraybackslash}X@{}}
\toprule
\textbf{Pengaturan} & \textbf{Nilai yang lazim} & \textbf{Pengaruhnya} \\
\midrule
\emph{Sprite Mode} &
  \emph{Single} atau \emph{Multiple} &
  Menentukan apakah berkas berisi satu gambar atau banyak gambar tersusun
  sebagai \emph{sprite sheet}. \\
\addlinespace
\emph{Pixels Per Unit} &
  Misalnya 64 atau 100 &
  Menentukan berapa piksel setara satu satuan pada \emph{scene}; nilai inilah
  yang mengaitkan ukuran gambar dengan ukuran level. \\
\addlinespace
\emph{Filter Mode} &
  \emph{Point} &
  Menentukan cara gambar dihaluskan saat diperbesar; \emph{Point}
  mempertahankan ketajaman piksel. \\
\addlinespace
\emph{Compression} &
  \emph{None} &
  Menentukan cara gambar dipadatkan saat program berjalan; \emph{None}
  menjaga mutu gambar. \\
\bottomrule
\end{tabularx}}
\vspace{0.10cm}
\begin{kotak}
Gambar 128 kali 128 piksel dengan \emph{Pixels Per Unit} 100 tingginya
$128 \div 100 = 1{,}28$ satuan. Agar tingginya 2 satuan, nilai yang diperlukan
$128 \div 2 = 64$. Samakan nilai itu untuk seluruh aset.
\end{kotak}
\end{frame}

\begin{frame}{\emph{Sprite Sheet} dan Pemotongannya}
\awalan{Sebuah \emph{sprite sheet} memuat beberapa gambar dalam satu berkas,
dan pemotongannya diatur melalui \emph{Sprite Editor}.}
\begin{itemize}
  \item Cara itu dipakai untuk gambar yang berubah-ubah, misalnya gerak
        berjalan yang dipecah menjadi beberapa bingkai.
  \item \emph{Sprite Mode} diatur menjadi \emph{Multiple}, lalu bingkainya
        dipotong pada \emph{Sprite Editor}.
  \item Tentukan ukuran setiap bingkai secara seragam, misalnya 64 kali 64
        piksel, sehingga pemotongannya dapat dilakukan otomatis.
  \item Ukuran yang tidak seragam menuntut pemotongan satu per satu; itu
        melelahkan dan mudah salah.
\end{itemize}
\vspace{0.12cm}
\begin{peringatan}
Jangan mengubah ukuran berkas gambar di luar Unity, misalnya dengan program
penyunting gambar. Apabila berkasnya berubah, Unity mengolahnya ulang dan
hasilnya untuk seluruh benda yang memakainya ikut berubah. Ubahlah melalui
\emph{Pixels Per Unit} agar perubahannya terkendali dan dapat dikembalikan.
\end{peringatan}
\end{frame}

% ========================================================= bagian 2
\section{Animator dan Peralihan Animasi}
\sectionframe{02}{Animator dan Peralihan Animasi}
  {Mesin keadaan yang mengurus gambar yang ditampilkan}

\begin{frame}{Mesin Keadaan untuk Animasi}
\awalan{Komponen \emph{Animator} adalah mesin keadaan yang khusus mengurus
animasi. Yang berpindah bukan perilaku benda, melainkan gambar yang
ditampilkan.}
\centering
\adjustbox{max width=0.90\linewidth,max height=0.52\textheight}{%
  % chapters/14/figures/animator-dan-canvas.tikz
% Susunan dua hal yang dibahas pada bab ini: mesin keadaan animasi, dan
% antarmuka pengguna beserta susunan bendanya.
\begin{tikzpicture}[
  keadaan/.style={
    rectangle, rounded corners=8pt, draw=ubPrimary, thick, fill=ubAlt,
    align=center, text width=18mm, inner sep=4pt, minimum height=10mm,
    font=\footnotesize
  },
  panah/.style={-{Stealth[length=2.2mm]}, thick, draw=ubAccent},
  label/.style={font=\scriptsize, align=center},
  judul/.style={font=\footnotesize\bfseries, ubPrimary, align=center}
]

% ---------- Panel A: mesin keadaan animasi ----------
\begin{scope}
  \node[judul, anchor=south] at (2.5,1.5) {mesin keadaan animasi};

  \node[keadaan] (diam) at (0,0) {diam\\(idle)};
  \node[keadaan] (lari) at (5.0,0) {berlari\\(run)};
  \node[keadaan] (lompat) at (2.5,-2.0) {melompat\\(jump)};

  \draw[panah] (diam) -- node[label, above, yshift=2pt]
    {kecepatan $>$ 0,1} (lari);
  \draw[panah] (lari) to[bend left=20] node[label, above, yshift=1pt]
    {kecepatan $\approx$ 0} (diam);
  \draw[panah] (lari) -- node[label, right, xshift=3pt]
    {tanah = salah} (lompat);
  \draw[panah] (lompat) -- node[label, left, xshift=-3pt]
    {tanah = benar} (diam);
\end{scope}

% ---------- Panel B: susunan antarmuka ----------
\begin{scope}[xshift=9.4cm]
  \node[judul, anchor=south] at (2.0,1.5) {susunan antarmuka pengguna};

  \node[keadaan, text width=30mm] (canvas) at (2.0,0.6) {Canvas};
  \node[keadaan, text width=16mm] (skor) at (0.3,-1.0) {Teks Skor};
  \node[keadaan, text width=16mm] (nyawa) at (2.0,-1.0) {Teks Nyawa};
  \node[keadaan, text width=16mm] (ikon) at (3.7,-1.0) {Ikon Nyawa};

  \draw[panah] (canvas) -- (skor);
  \draw[panah] (canvas) -- (nyawa);
  \draw[panah] (canvas) -- (ikon);

  \node[label, anchor=north, text width=50mm] at (2.0,-1.7)
    {Seluruh elemen antarmuka ditempatkan di bawah satu \istilah{Canvas},
     sehingga susunannya tetap sama pada setiap ukuran layar.};
\end{scope}
\end{tikzpicture}
}
\vspace{0.10cm}
\begin{columns}[T,onlytextwidth]
  \begin{column}{0.48\textwidth}
    {\scriptsize Keadaan pada \emph{Animator} disebut \emph{animation clip}:
     rangkaian gambar yang diputar berulang.}
  \end{column}
  \begin{column}{0.48\textwidth}
    {\scriptsize Transisinya diatur melalui \emph{parameter}, yaitu nilai
     yang ditetapkan oleh kode Anda.}
  \end{column}
\end{columns}
\end{frame}

\begin{frame}{Tiga Jenis Parameter}
\awalan{Pilihan yang tepat di antara ketiganya menghemat pekerjaan.}
\begin{itemize}
  \item \textbf{Angka pecahan} (\emph{float}) dipakai untuk nilai yang berubah
        halus, misalnya besar kecepatan.
  \item \textbf{Nilai benar atau salah} (\emph{bool}) dipakai untuk keadaan
        yang hanya memiliki dua kemungkinan, misalnya sedang menyentuh tanah
        atau tidak.
  \item \textbf{Pemicu} (\emph{trigger}) dipakai untuk kejadian yang terjadi
        sekali dan tidak memiliki lama, misalnya lompatan atau menerima
        kerusakan.
\end{itemize}
\vspace{0.12cm}
\begin{peringatan}
Sebuah lompatan sebaiknya memakai \emph{trigger}, bukan \emph{bool} yang harus
dikembalikan ke nilai semula oleh kode, sebab \emph{trigger} dikembalikan
sendiri oleh mesin setelah dipakai.
\end{peringatan}
\begin{kotak}
Nama parameter pada kode harus sama persis dengan nama pada jendela
\emph{Animator}, termasuk besar-kecil hurufnya. Kesamaan itu tidak diperiksa
penyunting.
\end{kotak}
\end{frame}

\begin{frame}[fragile]{Menghubungkan Animasi dengan Keadaan Pemain}
\lstinputlisting[language=CSharp,firstline=32,lastline=46]
  {../buku_ajar/code/unity/ch14_kontrol_animasi.cs}
\kodejelas{Skrip tidak memutar animasi apa pun. Ia hanya menetapkan nilai
parameter, sedangkan peralihannya dikerjakan mesin menurut transisi yang Anda
gambar pada jendela \emph{Animator}.}
\begin{itemize}
  \item Nama parameter disimpan pada peubah tetap melalui
        \texttt{Animator.StringToHash}, sehingga kesalahan penamaan dicegah
        sejak awal.
\end{itemize}
\end{frame}

% ========================================================= bagian 3
\section{Canvas: Skor dan Nyawa}
\sectionframe{03}{Canvas: Skor dan Nyawa}
  {Wadah tersendiri untuk antarmuka pengguna}

\begin{frame}{Mengapa Antarmuka Memerlukan Wadah Tersendiri}
\awalan{Elemen antarmuka berbeda sifatnya dari benda permainan.}
\begin{itemize}
  \item Ia tidak berada pada ruang permainan, tidak terpengaruh kamera, dan
        tidak bergerak ketika pemain berpindah.
  \item Karena itu ia ditempatkan pada wadah tersendiri yang disebut
        \emph{Canvas}.
  \item Seluruh elemen berada di bawah satu \emph{Canvas}, dan itulah yang
        membuat susunannya tetap sama pada setiap ukuran layar.
  \item Antarmuka permainan semacam ini lazim disebut \emph{HUD}
        (\emph{heads-up display}): skor, nyawa, dan keterangan yang melekat
        pada layar.
\end{itemize}
\vspace{0.10cm}
{\fontsize{7}{8.5}\selectfont\bfseries\color{ubAccent}TIGA ELEMEN YANG DIPAKAI}\par
\vspace{0.04cm}
\begin{itemize}
  \item \textbf{\emph{Text}} --- menampilkan tulisan, misalnya skor dan jumlah
        nyawa; ukuran dan jenis hurufnya diatur pada \emph{Inspector}.
  \item \textbf{\emph{Image}} --- menampilkan gambar, misalnya ikon nyawa atau
        bingkai antarmuka.
  \item \textbf{\emph{Panel}} --- wadah untuk mengelompokkan elemen, sehingga
        susunannya dapat dipindahkan bersama-sama.
\end{itemize}
\end{frame}

\begin{frame}[fragile]{Memperbarui Teks dan Ikon Nyawa}
\lstinputlisting[language=CSharp,firstline=46,lastline=59]
  {../buku_ajar/code/unity/ch14_tampilan_hud.cs}
\kodejelas{\texttt{PerbaruiTampilan} hanya menyalin keadaan ke tampilan: teks
skor dan nyawa diperbarui, lalu ikon nyawa dinyalakan menurut sisa nyawa.}
\begin{itemize}
  \item \texttt{Mathf.Max(0, lives - nilai)} menjaga nyawa tidak bernilai
        negatif, sebab angka negatif pada antarmuka terlihat janggal.
\end{itemize}
\end{frame}

\begin{frame}{Antarmuka Hanya Menampilkan, Bukan Menyimpan}
\awalan{Antarmuka tidak menyimpan keadaan permainan; ia hanya menampilkannya.}
\begin{itemize}
  \item Keadaan yang sebenarnya tetap disimpan pada kode Anda, dan antarmuka
        diperbarui ketika keadaannya berubah.
  \item Apabila skor disimpan pada tulisan di layar saja, keadaan permainan
        menjadi bergantung pada tampilannya, dan penyimpanan keadaan menjadi
        sulit.
  \item Aturannya sama dengan yang telah Anda pelajari pada Bab III: keadaan
        disimpan pada satu tempat yang jelas, sedangkan tampilan hanya
        mengikuti.
\end{itemize}
\vspace{0.12cm}
\begin{peringatan}
Rujukan ke elemen antarmuka dihubungkan melalui \emph{Inspector}, bukan melalui
kode. Apabila rujukan itu belum dihubungkan, program tetap berjalan tetapi
antarmukanya tidak diperbarui; kesalahannya muncul sebagai pesan peringatan
pada panel \emph{Console}. Periksa panel itu setiap kali antarmuka tidak
berubah seperti yang diharapkan.
\end{peringatan}
\end{frame}

% ========================================================= bagian 4
\section{Mengganti Greybox dengan Aset}
\sectionframe{04}{Mengganti Greybox dengan Aset}
  {Urutan tetap agar tidak ada yang terlewat}

\begin{frame}{Urutan Penggantian}
\awalan{Penggantian balok abu-abu dengan aset sebaiknya dilakukan menurut
urutan yang tetap.}
\begin{enumerate}
  \item \textbf{Pastikan mekaniknya sudah selesai.} Seluruh ukuran telah
        ditetapkan, dan levelnya sudah dapat dimainkan dengan lancar.
  \item \textbf{Impor seluruh asetnya lebih dahulu}, lalu atur
        \emph{importer setting}-nya satu per satu. Mengatur pengaturan sekaligus
        pada awal pekerjaan lebih cepat daripada mengaturnya sambil mengganti
        benda.
  \item \textbf{Ganti benda satu per satu}, mulai dari yang paling terlihat
        pemain: pemain, lalu rintangan, lalu latar. Setelah setiap penggantian,
        jalankan \emph{scene}-nya dan pastikan mekaniknya tidak berubah.
  \item \textbf{Periksa bentuk batas bendanya.} Mengganti gambar benda tidak
        mengubah \emph{Collider}-nya secara otomatis.
\end{enumerate}
\vspace{0.10cm}
\begin{kotak}
Langkah keempat paling sering terlewat. Gejalanya khas: benda yang sudah
bergambar rapi tetap bertabrakan seolah-olah masih berupa balok sebelumnya.
\end{kotak}
\end{frame}

\begin{frame}{Membandingkan Sebelum dan Sesudah}
\awalan{Perbandingan sebelum dan sesudah penggantian aset merupakan bagian dari
luaran praktikum pada RPS.}
\begin{itemize}
  \item \textbf{Memastikan mekaniknya tidak berubah.} Apabila pemain tiba-tiba
        tidak dapat melompati celah yang sebelumnya terjangkau, kemungkinan
        besar bentuk batas bendanya berubah bersama gambarnya.
  \item \textbf{Memeriksa keputusan perancangan.} Sering kali aset yang
        sesungguhnya menuntut penyesuaian pada levelnya: celah yang tadinya
        wajar menjadi terlalu sempit, atau warna latar bertabrakan dengan
        warna pemain.
  \item Temuan semacam itu wajar, dan justru merupakan alasan mengapa level
        disusun lebih dahulu secara \emph{greybox}.
\end{itemize}
\vspace{0.10cm}
\begin{kiat}
Simpan tangkapan layar sebelum dan sesudah penggantian aset, lalu letakkan
keduanya berdampingan pada laporan praktikum. Laporan yang memuat bukti
perubahan lebih meyakinkan daripada yang hanya menceritakannya.
\end{kiat}
\end{frame}

\begin{frame}{Bentuk Batas Tidak Ikut Berubah}
\awalan{Mengganti gambar benda tidak mengubah \emph{Collider}-nya secara
otomatis.}
\begin{itemize}
  \item Apabila ukuran gambarnya berbeda dari balok sebelumnya, bentuk batasnya
        perlu disesuaikan.
  \item Periksa letak dan ukuran \emph{Collider} pada \emph{Inspector}, lalu
        bandingkan dengan bentuk gambar yang baru.
  \item Jalankan \emph{scene}-nya setelah setiap penyesuaian; satu perubahan
        bentuk batas dapat mengubah lompatan, celah, dan daerah aman sekaligus.
  \item Apabila Anda memakai aset buatan orang lain, catatlah sumbernya beserta
        lisensinya sejak awal, ketika asetnya baru diimpor.
\end{itemize}
\vspace{0.12cm}
\begin{peringatan}
Mencari kembali asal-usul ratusan berkas pada akhir semester jauh lebih
melelahkan daripada mencatatnya satu per satu ketika diimpor.
\end{peringatan}
\end{frame}

% ========================================================= bagian 5
\section{Program Animasi dan Antarmuka}
\sectionframe{05}{Program Animasi dan Antarmuka}
  {Kode menentukan keadaan, mesin mengerjakan peralihan}

\begin{frame}[fragile]{Fungsi Publik yang Dipanggil Skrip Lain}
\lstinputlisting[language=CSharp,firstline=29,lastline=34]
  {../buku_ajar/code/unity/ch14_tampilan_hud.cs}
\kodejelas{\texttt{TambahSkor} adalah fungsi publik yang dipanggil dari skrip
lain, misalnya ketika pemain mengambil benda. Skor disimpan pada kode, lalu
tampilannya diperbarui.}
\begin{itemize}
  \item Pembagian pekerjaan pada animasi sama dengan pembagian pada Bab IX:
        kode menentukan keadaan, sedangkan peralihannya ditentukan aturan yang
        telah Anda susun.
  \item Arah hadap gambar dibalik dengan mengubah \texttt{transform.localScale},
        bukan memutar bendanya; memutar gambar dua dimensi 180 derajat akan
        membuatnya terbalik atas-bawah.
\end{itemize}
\vspace{0.10cm}
\begin{kotak}
\texttt{enabled = false} dipakai ketika rujukan antarmuka belum dihubungkan,
sehingga kesalahan yang sama tidak terulang pada setiap bingkai. Menghentikan
pekerjaan yang tidak dapat dilanjutkan lebih baik daripada membiarkannya
berjalan.
\end{kotak}
\end{frame}

\begin{frame}{Menelusuri Animasi yang Tidak Berganti}
\awalan{Gejalanya khas: program berjalan tanpa pesan kesalahan apa pun, tetapi
animasinya tidak pernah berganti.}
\begin{enumerate}
  \item \textbf{Periksa jendela \emph{Animator}.} Apakah kedua keadaannya ada,
        dan apakah transisinya tergambar?
  \item \textbf{Periksa nilai parameter ketika \emph{scene} berjalan.} Jendela
        itu menampilkan nilai setiap parameter pada saat itu.
  \item \textbf{Periksa kode yang menetapkan nilainya.} Nama \texttt{speed}
        pada kode berbeda dari \texttt{Speed} pada \emph{Animator}; perbedaan
        besar-kecil huruf tidak diperiksa.
  \item \textbf{Perbaiki, lalu jalankan kembali.} Animasinya berganti
        sebagaimana mestinya.
\end{enumerate}
\vspace{0.10cm}
\begin{kotak}
Jendela \emph{Animator} yang menampilkan nilai parameter saat berjalan adalah
alat pemeriksa yang paling berguna di sini. Kesalahan penamaan tidak
menimbulkan pesan apa pun.
\end{kotak}
\end{frame}

\begin{frame}{Percobaan yang Disarankan}
{\scriptsize
\renewcommand{\arraystretch}{1.12}
\begin{tabularx}{\textwidth}{@{}>{\raggedright\arraybackslash}p{6.4cm}
    >{\raggedright\arraybackslash}X@{}}
\toprule
\textbf{Yang diubah dan caranya} & \textbf{Yang sebaiknya diamati} \\
\midrule
\emph{Pixels Per Unit} pada aset pemain, dari 100 menjadi 50 &
  Gambar pemain menjadi dua kali lebih besar, sedangkan bentuk batasnya tidak
  ikut berubah. Sesuaikan bentuk batasnya sesudah percobaan. \\
\addlinespace
Nama parameter pada kode dibuat berbeda dari jendela \emph{Animator} &
  Animasi tidak pernah berganti, dan program tetap berjalan tanpa satu pun
  pesan kesalahan. Kembalikan sesudah percobaan. \\
\addlinespace
Nilai ambang kecepatan pada transisi dinaikkan menjadi 5 &
  Animasi berlari baru muncul ketika pemain benar-benar cepat. Perhatikan
  perbedaan rasa geraknya. \\
\addlinespace
Rujukan teks pada antarmuka diputus dari \emph{Inspector} &
  Pesan peringatan muncul pada panel \emph{Console}, dan antarmukanya tidak
  diperbarui. Kembalikan sesudah percobaan. \\
\bottomrule
\end{tabularx}}
\vspace{0.10cm}
\begin{peringatan}
Percobaan kedua paling penting untuk dicoba, sebab kesalahannya tidak
menimbulkan pesan apa pun. Peubah tetap untuk nama parameter mencegahnya sejak
awal.
\end{peringatan}
\end{frame}

% ========================================================= bagian 6
\section{Aset dan Hak Pakainya}
\sectionframe{06}{Aset dan Hak Pakainya}
  {Lisensi, atribusi, dan penutup bagian visual}

\begin{frame}{Tiga Jenis Lisensi yang Perlu Dikenal}
\awalan{Aset yang ditemukan di internet tidak selalu boleh dipakai. Periksa
lisensinya sebelum memakai aset apa pun.}
{\scriptsize
\renewcommand{\arraystretch}{1.14}
\begin{tabularx}{\textwidth}{@{}>{\raggedright\arraybackslash}p{3.3cm}
    >{\raggedright\arraybackslash}p{4.3cm}
    >{\raggedright\arraybackslash}X@{}}
\toprule
\textbf{Lisensi} & \textbf{Boleh dipakai untuk} & \textbf{Kewajiban} \\
\midrule
\emph{Public domain} atau \emph{CC0} &
  Segala keperluan, termasuk komersial. &
  Tidak ada. Cara paling bebas, dan paling aman untuk proyek mata kuliah. \\
\addlinespace
\emph{Creative Commons} dengan atribusi (CC~BY) &
  Segala keperluan, termasuk komersial. &
  Wajib menyebutkan pembuatnya dan lisensinya; untuk mata kuliah ini cukup
  pada laporan. \\
\addlinespace
\emph{Creative Commons} tanpa penggunaan komersial (CC~BY-NC) &
  Keperluan pendidikan, termasuk tugas mata kuliah. &
  Wajib menyebutkan pembuatnya, dan tidak boleh dipakai untuk produk yang
  dijual. \\
\bottomrule
\end{tabularx}}
\vspace{0.08cm}
{\footnotesize\color{ubMuted}\itshape Bacalah lisensi setiap aset pada halaman
asalnya; sebagian aset menyertakan berkas lisensi tersendiri.}
\end{frame}

\begin{frame}{Mencatat Sumber Aset}
\awalan{Catatan sumber aset sebaiknya dibuat sejak awal, dalam bentuk yang
sederhana.}
\begin{columns}[T,onlytextwidth]
  \begin{column}{0.48\textwidth}
    {\fontsize{7}{8.5}\selectfont\bfseries\color{ubAccent}EMPAT HAL YANG DICATAT}\par
    \vspace{0.06cm}
    \begin{itemize}
      \item Nama berkasnya.
      \item Sumbernya.
      \item Jenis lisensinya.
      \item Keperluan atribusinya.
    \end{itemize}
  \end{column}
  \begin{column}{0.48\textwidth}
    {\fontsize{7}{8.5}\selectfont\bfseries\color{ubAccent}GUNA CATATAN ITU}\par
    \vspace{0.06cm}
    \begin{itemize}
      \item Memenuhi kewajiban atribusi.
      \item Memudahkan mencari pengganti apabila diperlukan.
      \item Menjadi bukti kejujuran akademik pada laporan proyek akhir.
    \end{itemize}
  \end{column}
\end{columns}
\vspace{0.14cm}
\begin{peringatan}
Aset yang ditemukan melalui mesin pencari gambar belum tentu boleh dipakai.
Hak cipta berlaku sejak karya itu dibuat, walaupun tidak ada keterangan apa pun
pada berkasnya. Pakai hanya aset dari tempat yang menyatakan lisensinya dengan
jelas.
\end{peringatan}
\end{frame}

\begin{frame}{Menutup Bagian Visual dari CPMK 4}
\awalan{Bab ini melengkapi bagian pertama CPMK 4, yaitu penyajian visual.}
{\scriptsize
\renewcommand{\arraystretch}{1.12}
\begin{tabularx}{\textwidth}{@{}>{\raggedright\arraybackslash}p{1.7cm}
    >{\raggedright\arraybackslash}p{5.0cm}
    >{\raggedright\arraybackslash}X@{}}
\toprule
\textbf{Bab} & \textbf{Yang ditambahkan} & \textbf{Akibatnya} \\
\midrule
Bab XIV & Aset visual, animasi, dan UI dasar &
  \emph{Placeholder} \emph{greybox} diganti aset sesungguhnya; pemain diam dan
  berlari; \emph{Canvas} menampilkan skor dan nyawa. \\
\addlinespace
Bab XV & Bunyi dan umpan balik interaksi &
  Peristiwa permainan---lompatan, pengambilan benda, kalah---memperoleh
  tanggapan bunyi. \\
\addlinespace
Bab XVI & Pengujian prototipe dan presentasi &
  Aset yang telah dipasang dinilai mutunya oleh pemain lain, dan hasilnya
  menjadi dasar perbaikan. \\
\bottomrule
\end{tabularx}}
\vspace{0.12cm}
\begin{kotak}
Yang tersisa adalah penyajian bunyi, dan pengujian atas keduanya secara
menyeluruh. Keduanya berpijak pada \emph{scene} yang baru saja Anda lengkapi.
\end{kotak}
\end{frame}

% --------------------------------------------------------- aktivitas
\begin{frame}{Aktivitas di Kelas \textperiodcentered\ Praktikum 11}
\setbeamerfont{itemize/enumerate body}{size=\scriptsize}
\awalan{Pengalaman belajar minggu 13 pada RPS, dikerjakan berurutan di kelas
lalu diselesaikan di rumah.}
\begin{enumerate}
  \item Impor sekurang-kurangnya satu aset pemain yang memuat animasi diam dan
        berlari, beserta aset latar dan rintangan. Catat sumber dan lisensinya
        sejak awal.
  \item Atur \emph{Pixels Per Unit} pada seluruh aset agar nilainya sama, lalu
        hitung tinggi pemain yang dikehendaki. Potong \emph{sprite sheet} menjadi
        bingkai berukuran seragam.
  \item Susun \emph{Animator}: keadaan diam dan berlari, parameter kecepatan dan
        keadaan menyentuh tanah, lalu gambar transisinya.
  \item Tempelkan \texttt{ch14\_kontrol\_animasi.cs}, hubungkan rujukan
        \emph{Animator} dan \texttt{Rigidbody} pada \emph{Inspector}, lalu uji
        pergantian animasinya.
  \item Buat \emph{Canvas} berisi teks skor dan nyawa, tempelkan
        \texttt{ch14\_tampilan\_hud.cs}, hubungkan rujukan teksnya, lalu panggil
        \texttt{TambahSkor} ketika pemain mengambil benda.
  \item Simpan tangkapan layar sebelum dan sesudah penggantian aset, periksa
        bentuk batasnya, tulis catatan praktikum, lalu simpan ke repositori Git.
\end{enumerate}
\vspace{0.06cm}
{\footnotesize\color{ubAccent}\itshape Rincian lengkap: Praktikum 11 pada
Bab XIV buku ajar. Kumpulkan melalui repositori Git pribadi.}
\end{frame}

% --------------------------------------------------------- rangkuman
\begin{frame}{Rangkuman}
\begin{itemize}
  \item Aset diolah ketika diimpor, bukan ketika program dijalankan. Mengubah
        \emph{importer setting} mengubah seluruh benda yang memakai aset itu
        sekaligus.
  \item Empat pengaturan impor yang paling berpengaruh: \emph{Sprite Mode},
        \emph{Pixels Per Unit}, \emph{Filter Mode}, dan \emph{Compression}.
        Nilai \emph{Pixels Per Unit} sebaiknya sama untuk seluruh aset.
  \item \emph{Sprite sheet} dipotong menjadi bingkai berukuran seragam melalui
        \emph{Sprite Editor}; ukuran berkas tidak boleh diubah di luar Unity.
  \item \emph{Animator} adalah mesin keadaan animasi. Keadaannya berupa
        \emph{animation clip}, dan transisinya diatur melalui parameter
        \emph{float}, \emph{bool}, atau \emph{trigger}.
  \item Nama parameter harus sama persis dengan yang ada pada jendela
        \emph{Animator}; kesalahannya tidak menimbulkan pesan apa pun.
  \item Antarmuka pada \emph{Canvas} hanya menampilkan keadaan yang disimpan
        pada kode. Mengganti aset menuntut pemeriksaan bentuk batas benda.
\end{itemize}
\end{frame}

% --------------------------------------------------------- cek
\begin{frame}{Cek Pemahaman}
\begin{enumerate}
  \item Sebuah gambar berukuran 192 kali 192 piksel diimpor dengan
        \emph{Pixels Per Unit} bernilai 100. Hitunglah tingginya pada
        \emph{scene}, lalu hitunglah nilai yang diperlukan agar tingginya
        menjadi 2 satuan.
  \item Kode menetapkan nilai parameter animasi, tetapi animasinya tidak
        berganti dan tidak ada pesan kesalahan. Sebutkan dua kemungkinan
        penyebabnya.
  \item Mengapa keadaan permainan disimpan pada kode, bukan pada tampilan
        antarmukanya?
\end{enumerate}
\vspace{0.35cm}
{\footnotesize\color{ubMuted}\itshape Jawab tanpa membuka catatan, lalu
bandingkan dengan pasangan di sebelah Anda.}
\end{frame}

% --------------------------------------------------------- penutup
\penutup{Tugas Praktikum 11}
  {Materi 13 \textperiodcentered\ Integrasi Audio dan \emph{Trigger} Interaksi:
   \emph{AudioClip} dan \emph{AudioSource} mengiringi BGM serta SFX, lalu
   \emph{trigger} pada \emph{pickup}, tabrakan, dan kondisi menang atau kalah
   memberi umpan balik atas peristiwa permainan.}

\end{document}
